\exercice{Orage}
\enonce{
Au cours d'un orage, un éclair peut être assimilé à un conduit cylindrique rectiligne de rayon $a=\SI{10}{cm}$ et parcouru par un courant $I=\SI{1e5}{A}$ de densité volumique de courant uniforme. On supposera dans un premier temps le courant ascendant. Les effets de bord seront négligés.}

\begin{enumerate}
  \question{Faire un schéma et expliqué pourquoi les électrons de ce courant sont soumis à une force magnétique de Lorentz. Quel est son sens ?}
    [Quelle est la direction de la vitesse des électrons ? Quelle est la direction du champ magnétique créé par le courant ?]
  \question{Exprimer le vecteur densité volumique de courant $\vv{j}$.}
    [Relier le courant à $\vv{j}$. Utiliser le fait que $\vv{j}$ est uniforme dans le cylindre.]
  \question{Déterminer $\vv{B}$ au niveau du bord du conduit et exprimer la norme de cette force magnétique par unité de volume en fonction de $I$ et $a$.}
    [Effectuer les 4 étapes : analyse des invariances, analyse des symétries, choix de la courbe d'Ampère, théorème d'Ampère.]
    [On demande $\vv{B}$ seulement au bord du cylindre (en $r=a$).]
  \question{Le sens de la force change-t-il si le courant est descendant ?}
    [Quel serait alors le sens de la vitesse des électrons ? Quel serait alors le sens du champ magnétique ?]
  \question{Faire l'application numérique de cette force et la comparer au poids volumique de l'air. Pourquoi les éclairs causent-ils le tonnerre\footnote{L'éclair est le résulultat visible du passage du courant tandisque le tonnerre est le son produit} ?}
    [Relier le poids volumique de l'air à la masse volumique de l'air.]
    [La masse volumique de l'air est $\rho_\text{air}=1kg.m^{-3}$.]
    [Que peut-il se produire si l'air est subitement soumis à une force très intense ?]
\end{enumerate}

\exercice{Bobine torique}
\enonce{Une bobine est constituée par un fil conducteur bobiné en spires jointives sur un tore circulaire à section carrée de côté $a$, de rayon moyen $R$. On désigne par $n$ le nombre total de spires et par $I$ le courant qui les parcourt. On note $(Oz)$ son axe de symétrie.}

\enonce{On s'intéresse au champ magnétique à l'intérieur du tore.}

\begin{enumerate}
  \question{Faire un schéma du système.}
  \question{Montrer que le champ magnétique qui règne en un point $M(r,z)$ quelconque à l'intérieur du tore peut s'exprimer sous la forme  $B=\frac{\mu_0nI}{2\pi r}$.}
    [Effectuer les 4 étapes : analyse des invariances, analyse des symétries, choix de la courbe d'Ampère, théorème d'Ampère.]
  \question{Déterminer le flux $\Phi$ du champ magnétique à travers la surface d'\textbf{une} spire dont la normale est orientée dans le sens du champ.}
    [Représenter la surface sur le schéma.]
    [Selon quelles variables faut-il intégrer ? Entre quelles bornes ?]
    [Quelle est l'expression d'un élément de surfaces orienté selon $\vv{e_\theta}$]
\end{enumerate}

\exercice{Câble coaxial}
\enonce{On considère un câble coaxial cylindrique de longueur supposée infinie, constitué d'un conducteur central plein de rayon $R_1$, parcouru par un courant uniforme d'intensité $I$ et d'un conducteur périphérique évidé, de rayon intérieur $R_2$, de rayon extérieur $R_3$ avec $R_1<R_2<R_3$ et parcouru par un courant uniforme également d'intensité $I$ mais circulant en sens inverse par rapport au courant du conducteur central.}

\enonce{On notera $\vv{u_z}$ le vecteur unitaire de l'axe commun des deux conducteurs. Soit un point $M$ à une distance $r$ de l'axe du câble.}

\begin{enumerate}
  \question{Montrer que le champ magnétique $\vv{B}$ créé au point $M$ est orienté selon $\vv{e_\theta}$.}
  \question{Montrer qu'il peut se mettre sous la forme $\vv{B}=B(r)\vv{u_\theta}$.}
  \question{Préciser alors la forme des lignes de champ.}
    [Les lignes de champ sont en tout point colinéaires à $\vv{B}$.]
  \question{Montrer que le champ magnétique créé au point $M$ est nul si $r>R_3$.}
    [Quel est le courant enlacé dans ce cas ?]
  \question{Calculer les densités de courant $\vv{j_1}$ et $\vv{j_2}$, respectivement dans le conducteur central et du conducteur périphérique en fonction des courants $I$ et $-I$ et des rayons $R_1$, $R_2$ et $R_3$.}
    [Quelle est l'aire du disque de rayon $R_1$ ? Quelle est l'aire du disque percé situé entre $R_2$ et $R_3$ ?]
    [Relier le courant électrique et la densité volumique de courant $j$.]
  \question{En appliquant le théorème d'Ampère à un contour $\mathcal{C}$ que l'on précisera, donner l'expression de la composante $B(r)$ du champ magnétique créé au point $M$ en fonction de $\mu_0$, $I$, $r$, $R_1$, $R_2$, et $R_3$ dans chacun des cas suivants : $r<R_1$ ; $R_1<r<R_2$ et $R_2<r<R_3$.}
    [La courbe d'Ampère doit être choisie de sorte que $\vv{dl}$ soit colinéaire à $\vv{B}$ afin de simplifier les calculs.]
  \question{Tracer l'allure du graphe de $B(r)$.}
\end{enumerate}

\exercice{Mesure du champ magnétique terrestre}
\enonce{Pour mesurer la composante horizontale du champ magnétique terrestre, on uilise une boussole. On place un solénoïde autour de la boussole de sorte que l'axe du solénoïde soit horizontal et orthogonal à l'aiguille de la boussole. Lorsqu'on fait circuler un courant dans le solénoïde, l'aiguille tourne d'un angle $\alpha=\SI{40}{\deg}$.}

\enonce{Données : rayon du solénoïde $R=\SI{5}{cm}$, nombre de spires $N=100$, longueur $l=\SI{30}{cm}$, intensité du courant $I=\SI{1}{mA}$.}

\begin{enumerate}
  \question{Calculer la composante normale du champ magnétique terrestre.}
    [Faire un schéma.]
    [Exprimer le champ magnétique créé par le solénoïde, puis le champ total.]
\end{enumerate}

% Exo avec simulation

% J'explique à mon oncle